\documentclass[10pt,twocolumn]{article}
\usepackage[margin=1.9cm,top=2.2cm]{geometry}
\usepackage{graphicx}
\usepackage{amsmath,amssymb}
\usepackage{microtype}
\usepackage{lmodern}
\usepackage{xcolor}
\usepackage{booktabs}
\usepackage{caption}
\usepackage[colorlinks=true,linkcolor=blue!50!black,citecolor=blue!50!black,urlcolor=blue!50!black]{hyperref}
\captionsetup{font=small,labelfont=bf,skip=6pt}
\setlength{\parskip}{0pt}
\usepackage{titlesec}
\setcounter{topnumber}{4}
\setcounter{totalnumber}{6}
\renewcommand{\topfraction}{0.95}
\renewcommand{\textfraction}{0.03}
\renewcommand{\floatpagefraction}{0.85}
\titleformat{\section}{\large\bfseries}{\thesection}{0.6em}{}
\titleformat{\subsection}{\normalsize\bfseries}{\thesubsection}{0.6em}{}
\titlespacing{\section}{0pt}{10pt}{5pt}
\titlespacing{\subsection}{0pt}{8pt}{3pt}

\title{\Large\bfseries From pairwise electrocoalescence to optical clearing:\\
a predictive dimensionless framework for electrostatic\\ demulsification of nanoemulsions}
\author{Sreedath Panat$^{1,\ast}$, Raj Dandekar$^{1}$, Rajat Dandekar$^{1}$\\[2pt]
\small $^{1}$Vizuara AI Labs\\
\small $^{\ast}$Corresponding author.}
\date{}

\begin{document}
\twocolumn[
\maketitle
\begin{center}\begin{minipage}{0.94\textwidth}\small
\noindent\textbf{Abstract.}
The pair-scale physics of electrocoalescence is textbook material: water
droplets in oil polarize, attract, chain, and merge. Yet the observable that
decides whether an electrostatic demulsifier works, the time for the emulsion
to \emph{optically clear}, cannot be obtained from any single-pair argument,
and has so far been measured rather than predicted. Here we construct a
predictive chain from pair electrostatics to the operational 95\%-transmittance
endpoint for water-in-oil nanoemulsions, using only pre-experiment setup
information: material properties, cell geometry, device characterization, the
initial droplet-size distribution, and a bound on the surfactant film barrier
extracted from the emulsion's zero-field shelf life. Three independent
computational routes, a closed-form staged cascade, a polydisperse population
balance with a Mie-scattering endpoint, and a stochastic super-droplet
simulation steered by a reinforcement-learning agent, predict, without
calibration to any kinetic measurement, clearing in 90--125 minutes for a
175\,nm, 2\% emulsion at 7.5\,kV/cm; the measured value, not consulted
during model construction, is $\sim$60 minutes. The
framework identifies two well-separated clocks: bulk water separation
(minutes) and optical clearing (hours), the latter controlled by a
sub-100-nm fines tail. A single dimensionless group, the electrocoalescence
number $\Lambda=\pi\varepsilon_0\varepsilon_{\mathrm{oil}}K^2E^2a^3/k_BT$,
gates the entire cascade: published voltage-sweep data collapse onto
$t\cdot\Lambda=23.5$\,min (5\% RMS scatter), with the go/no-go threshold
$\Lambda^{*}\!\approx\!0.8$ predicted in advance. The framework further
explains, without fitting, why saline droplets demulsify more slowly than
fresh water. Code and data are openly available.
\end{minipage}\end{center}
\vspace{14pt}
]

\section{Introduction}

Separating emulsified water from oil is a rate-limiting operation across crude
desalting, produced-water treatment, and waste-oil recycling
\cite{kokal2005,zolfaghari2016}, and it becomes hardest exactly where it
matters most: at nanoscale droplet sizes, where gravity is powerless and
surfactant films render the dispersion metastable for
months~\cite{eow2001,tadros2009,mason2006,gupta2016}. Electrocoalescence,
driving droplet merger with an electric field, is the classical
remedy~\cite{atten1993,eow2001,mhatre2015,less2012}, but
conventional immersed-electrode devices are capped near 1\,kV/cm by
shorting through droplet chains. A recently demonstrated non-Laplacian,
air-gap corona architecture removes that cap and sustains fields up to
$\sim$7.5\,kV/cm across the emulsion, demulsifying 150--800\,nm water-in-oil
emulsions without chemical additives~\cite{panat2025}.

The microphysics underlying such devices is well established. A conducting
aqueous droplet in oil acquires an induced dipole; aligned dipoles attract
and chain, a fibration known since Winslow's early electrorheology
\cite{winslow1949,parthasarathy1996}, and press out the intervening oil
film until rupture and merger follow~\cite{atten1993,zhang1995,jones1995,
lathamroxburgh1966}. What is \emph{not} established is
the bridge from this pair-scale picture to the quantity an operator actually
observes: the time for a milky emulsion to become transparent. The two are
separated by seven orders of magnitude in droplet count, by a coalescence
cascade spanning three decades in radius, by gravity, and by the optics of
light scattering. Population-balance models of electrocoalescers have made
important progress on parts of this bridge, predicting droplet-size
evolution and water-removal efficiency in microscale
systems~\cite{sharma2023,anand2019,su2023,rahmani2023,khajeh2020}, but they
are calibrated against measured kinetics, address micron-scale emulsions,
and stop short of the observable an operator actually uses: the optical
clearing of the sample. For the nanoemulsion regime, where Brownian motion,
surfactant film barriers, and light scattering by sub-wavelength droplets
dominate, the clearing time has remained an empirical quantity: measured
and correlated with voltage, but not predicted from first principles.

This paper constructs and tests that bridge. Our contributions are fourfold.
First, a multiscale framework that converts setup information alone,
material properties, geometry, device current--voltage characterization, the
initial size distribution, and a film-barrier bound extracted from the
emulsion's zero-field shelf life, into a prediction of the \emph{optical}
clearing time with zero calibrated parameters; unlike prior
population-balance treatments~\cite{sharma2023,anand2019,su2023}, nothing
is fitted to kinetic data, the endpoint is evaluated through Mie optics
exactly as measured, and the regime is the Brownian-dominated nanoemulsion
limit. Second, a strict separation of two
observables that are usually conflated: the \emph{mass clock} (bulk water
resolved into a free layer, minutes at high field) and the \emph{optical
clock} (95\% relative transmittance, slower by a factor 20--50), the latter
controlled by a fines tail that no single-size argument can represent.
Third, a dimensionless organizing group: the electrocoalescence number
$\Lambda$, the induced-dipole pair energy at contact in units of $k_BT$.
$\Lambda$ is the droplet analogue of the dipolar coupling constant of
electrorheology~\cite{halsey1990}; we show that it gates feasibility with a
threshold predictable from shelf life, and that the published voltage-sweep
data of Ref.~\cite{panat2025} collapse onto $t\cdot\Lambda=\text{const}$ to
within $\pm5\%$. Fourth, a validation protocol we believe is itself useful: the prediction
was made without calibration to any kinetic measurement, by three
population-integration machineries with disjoint numerical failure modes
(analytic, population balance, and stochastic simulation driven by a
reinforcement-learning agent), with the measured kinetics quarantined until
all three numbers were frozen. We are explicit below about what this
protocol does and does not establish.

\begin{figure*}[!t]
\centering
\includegraphics[width=\textwidth]{figs/fig1_mechanism.png}
\caption{\textbf{Setup and the multiscale chain from pair electrostatics
to optical clearing.} (a)~The non-Laplacian batch configuration of
Ref.~\cite{panat2025}: a needle emitter array at $+V$ injects space charge
across an air gap of width $d$ onto a water-in-oil nanoemulsion of depth
$H$ above a grounded plate, sustaining fields up to
$E\approx7.5$\,kV/cm in the oil without shorting. (b)~Pair-scale
electrocoalescence: droplets polarize into induced dipoles, attract and
chain along the field, drain the intervening oil film against the
surfactant barrier $\Delta G\approx20$--$30\,k_BT$ (the rate-limiting
step), and merge. (c)~Population scale: a parallel cascade of $\sim$25
volume doublings, with gravitational sweep dominating above a micron,
delivers drops to Stokes settling; two clocks result, with bulk separation
in minutes and optical clearing delayed by a sub-100-nm residual fog. The
electrocoalescence number $\Lambda$ gates the entire chain.}
\label{fig:mechanism}
\end{figure*}

\section{Results}

\subsection{A framework built from setup facts alone}

We consider the batch configuration of Ref.~\cite{panat2025}: 60\,ml of
water-in-hexadecane emulsion (0.5\,M brine, 1\% Span~80, water cut
$\phi=2\%$, DLS mean diameter 175\,nm, lognormal width $\sigma_g=1.35$) in a
cubic cell of depth $H=3.37$\,cm, under a corona air-gap field of up to
$E=7.5$\,kV/cm in the oil. Material properties (viscosity $\mu$, densities,
permittivity $\varepsilon_{\mathrm{oil}}$, interfacial tension) are standard
literature values (Methods). Crucially, no measured demulsification kinetics
enter the model at any point; an older reduced-order model calibrated to the
measured timescale was deliberately quarantined.

One further setup fact does quantitative work. The emulsion shows no visible
phase separation for months without a field. Since the zero-field Brownian
encounter rate is known, stability over a shelf life $t_s$ bounds the
per-encounter coalescence probability, and hence the surfactant film barrier:
$\Delta G \gtrsim k_BT\ln(K_Bn_0t_s)\approx 19\,k_BT$ (instantiating
``several months'' as three; the criterion, fewer than one coalescence per
drop, is stricter than ``no visible separation'' and thus a conservative
lower bound). We carry
$\Delta G=20$--$30\,k_BT$ as the honest uncertainty band; Span-type
monolayers stiffened by electrolytes sit in this range~\cite{opawale1998}.

The pair physics is assembled from established elements: conductor-limit
induced dipoles~\cite{jones1995}, Brownian transport with
dipolar-capture enhancement (a Fuchs-type factor
$\sim\Lambda^{1/3}$)~\cite{fuchs1934,friedlander2000}, near-contact
amplification of the conductor--conductor attraction over the point-dipole
law~\cite{davis1964,zhang1995}, film-drainage barrier crossing, and, once
drops grow past a micron, sweep coagulation: settling drops collect smaller
partners from a capture tube whose radius is bounded jointly by energy
($U(r)\!=\!k_BT$) and by drift-versus-transit kinetics
(Methods)~\cite{pruppacher1997,friedlander2000}. Merges of conducting drops
in strong fields can also abort, ejecting charged daughter droplets; this
partial-coalescence channel is gated by the electric Bond number at contact
and by charge relaxation in the drop~\cite{ristenpart2009,bird2009,
mousavichoubeh2011}, and is retained because it carries the salinity
dependence discussed below.

\subsection{One number gates the cascade}

Every stage of the chain is a contest between the induced-dipole bond and
thermal noise. Dividing the aligned-pair contact energy by $k_BT$ defines
the electrocoalescence number,
\begin{equation}
\Lambda \;=\; \frac{\pi\,\varepsilon_0\,\varepsilon_{\mathrm{oil}}\,K^2E^2a^3}{k_BT},
\label{eq:lambda}
\end{equation}
with $E$ the field in the oil, $a$ the droplet radius, and $K\!\to\!1$ the
Clausius--Mossotti factor in the conductor limit. $\Lambda$ is the droplet
analogue of the dipolar coupling constant governing chain formation in
electrorheological fluids~\cite{halsey1990}. Coalescence, however, must also
defeat the film barrier. Two amplifications act at contact: a large drop
presses a small one with $\sim$8$\times$ the equal-pair energy (the
``scavenging rule'', Methods), and the near-contact conductor interaction
exceeds the point-dipole law by a factor $A_{nc}\!\approx\!2$--8
\cite{davis1964}. The go/no-go threshold is therefore
\begin{equation}
\Lambda^{*}\;\approx\;\frac{\Delta G/k_BT}{8\,A_{nc}}\;\approx\;0.8
\qquad(\text{band }0.3\text{--}2),
\label{eq:lambdastar}
\end{equation}
predicted before any kinetic data are consulted. The two amplifications
are not strictly independent (both are near-contact corrections to the
point-dipole law), so their product should be read as a heuristic
16--64$\times$ band whose uncertainty the threshold inherits; ignition at
$\Lambda\sim\Lambda^{*}$ additionally relies on the large-size tail of
the initial distribution, since an equal median pair at threshold still
faces a substantial residual barrier. Since
$\Lambda\!\propto\!E^2a^3$, the critical field steepens as droplets shrink,
$E_{\mathrm{crit}}\!\propto\!a^{-3/2}$ (Fig.~\ref{fig:map}): the compact
statement of why nanoemulsions defeat conventional electrocoalescers and
why the high-field air-gap architecture succeeds.

\begin{figure}[!t]
\centering
\includegraphics[width=\columnwidth]{figs/fig_map.png}
\caption{\textbf{The feasibility map.} The $\Lambda^{*}$ boundary
($E\!\sim\!a^{-3/2}$; orange band spans $\Lambda^{*}=0.3$--2) in
field-versus-size space, with operating points from
Ref.~\cite{panat2025}: filled, cleared; open, remained milky. The sweep's
cleared points bracket the boundary band, placing the operational
threshold at the band's lower edge; a 600\,nm emulsion clears at modest
fields because $\Lambda\!\propto\!a^{3}$.}
\label{fig:map}
\end{figure}

\subsection{Calibration-free prediction, three ways}

To make the prediction falsifiable we froze all model choices before
consulting the measured kinetics, and solved the same physics by three
machineries with disjoint numerical failure modes
(Fig.~\ref{fig:protocol}). We stress that this is a no-calibration rather
than a blind protocol: one author (S.P.) performed the source experiments
and knew the measured clearing time. The safeguards are procedural, namely
that no measured kinetic quantity enters the model at any point, that an
older reduced-order model calibrated to the measured timescale was
quarantined, and that all model choices and the three predicted numbers
were fixed, under version control, before any comparison to data. The
result therefore demonstrates that the framework requires no kinetic
calibration, not that the modellers were ignorant of the answer; a
hash-committed, pre-experiment version of this protocol would be stronger
and we recommend it for future systems.

\emph{Route 1, analytic cascade.} Closed-form staging: a Brownian gate
($\Lambda=5.2$ at the 87.5\,nm median radius and 7.5\,kV/cm: open;
$\Lambda=7.8$ at the 100\,nm mean radius used for the data collapse), barrier-limited ignition
(seconds), a growth ladder summing $\sim$25 volume-doubling times with
Brownian-capture and sweep transport, and a gravity endgame at the handoff
radius minimizing total time. Being single-sized, this route can only time the \emph{mass} clock:
5\,min (brine, barrier 20--25\,$k_BT$), rising to $\sim$40\,min at the
30\,$k_BT$ barrier edge.

\emph{Route 2, population balance.} A 123-bin sectional
Smoluchowski model~\cite{smoluchowski1917,kumar1996} with coagulation,
Stokes settling loss, a daughter-droplet source from partial coalescence,
and the endpoint evaluated exactly as defined operationally: relative
transmittance from full Mie scattering~\cite{bohren1983} reaching 95\% over
a 1\,cm path. Numerically, merges between comparable sizes are discrete
fixed-pivot events while accretion of much smaller partners is treated as
continuous growth (upwind advection in size space), which removes an
otherwise fatal stiffness; mass is conserved to 1\%.

\emph{Route 3, stochastic simulation with a learning agent.} An independent
super-droplet Monte Carlo implementation~\cite{shima2009} exposes only
operator-visible observables (turbidity, coarse mean size, free-water-layer
height) to a tabular reinforcement-learning agent~\cite{sutton2018} that
selects the field level every 30\,s under an arcing risk above 8\,kV/cm.
With reward equal to negative elapsed time, the learned value of the initial
state \emph{is} the predicted clearing time. The agent independently
rediscovers the experimental operating point, running at the arcing edge.

\begin{figure*}[!t]
\centering
\includegraphics[width=\textwidth]{figs/fig2_protocol.png}
\caption{\textbf{The blind-prediction protocol.} Setup facts (materials,
geometry, device characterization, size distribution, and a film-barrier
bound from zero-field shelf life) feed three independent routes: a
closed-form cascade, a population balance with Mie optics, and a
super-droplet Monte Carlo rig driven by a reinforcement-learning agent.
Measured kinetics remain quarantined until all three predictions are frozen.}
\label{fig:protocol}
\end{figure*}

The three routes agree with one another within 21\% (range over mean of
the central estimates) on the optical clock: population balance 110\,min
for brine and 90\,min for DI water (film-barrier band widens the brine
value to 78\,min--3.3\,h), fixed-field super-droplet Monte Carlo 102\,min (median over 20 random
seeds, observed range 81--273\,min), and learned value
estimate 125\,min. Throughout, the quoted prediction band 90--125\,min is
the range of central estimates across routes and water phases; it excludes
the film-barrier systematic, which is reported separately. Cross-route disagreements during development were diagnostic
rather than cosmetic: they exposed a mass leak at the size-grid boundary, an
event-rate undercount in rare sweep interactions, and the stiffness
pathology above, all of which a single-method study would have shipped
silently. A thirteen-point verification suite (dimensional checks against raw
constants, optical limits, grid-convergence, collapse statistics, and an
AST-based solver-independence audit) accompanies the code and returns a
nonzero exit status on any failure.

\subsection{Two clocks, one process}

The framework's central structural prediction is that ``separated'' is two
different claims (Fig.~\ref{fig:dynamics}). The cascade is parallel: all
pairs merge simultaneously, the population halves per generation, and only
$\log_2(4\times10^{7})\!\approx\!25$ generations connect 175\,nm to settling
sizes. Consequently bulk separation is fast: in the population balance, 97\% of
the water reaches the free layer within $\sim$2\,min at 7.5\,kV/cm (the
analytic route's central estimate is 5\,min). The transmittance endpoint, the standard optical route to emulsion
stability assessment~\cite{mengual1999,reddy1981}, is another matter. A residual fog holding under half a percent of the water,
droplets too small for even the amplified contact energy to open the film
barrier, scatters strongly enough to keep the sample milky, and it can no
longer coalesce in parallel: each straggler waits alone for a rare
collector. Logarithmic physics for the bulk; linear waiting for the tail.
The two clocks differ by a factor of 20--50, which is precisely why clearing
times have resisted single-droplet reasoning: the eye's endpoint is a
statement about the smallest decile of the initial size distribution, not
about the bulk. Staged kinetics with a slow residual-emulsion tail have been
reported in population-balance studies of microscale
electrocoalescers~\cite{sharma2023}; the present contribution is to make
the final stage quantitative through the scattering optics of the fines and
to show it is the operational endpoint's rate-limiting step.

\begin{figure*}[!t]
\centering
\includegraphics[width=0.9\textwidth]{figs/fig_cascade.png}
\caption{\textbf{The coalescence cascade (population balance,
7.5\,kV/cm, brine).} Volume-weighted size distributions march through
$\sim$25 volume doublings within seconds (blue curves, labelled by time);
by ten minutes (red dashed) the largest drops have settled out and the
suspended remainder is the residual fog plus the last settling drops.}
\label{fig:cascade}
\end{figure*}

\begin{figure*}[!t]
\centering
\includegraphics[width=0.9\textwidth]{figs/fig_transients.png}
\caption{\textbf{Two clocks (population balance, 7.5\,kV/cm).}
(a)~The optical clock: droplet optical depth decays toward the
95\%-transmittance threshold (dashed) in 90--110\,min, brine (red) slower
than DI water (blue). (b)~The mass clock: resolved water exceeds 97\%
within minutes for both phases.}
\label{fig:clocks}
\end{figure*}

\subsection{The data collapse onto one line}

Reference~\cite{panat2025} reports batch demulsification times over a
threefold voltage sweep, fitted there as $t\!\sim\!V^{-1.91}$. At fixed
droplet size the collapse below is therefore a dimensionless
\emph{restatement} of that fit rather than an independent measurement; its
new content is the value and tightness of the constant, the a-priori
threshold it terminates on, and the field-reading discrimination it
enables. Digitizing the seven published points and converting voltage to $\Lambda$
(field read as $V/d$; mean droplet radius $a=100$\,nm per the source's
caption; electrode gap $d=5.5$\,cm per its Methods) yields
\begin{equation}
t_{95}\cdot\Lambda \;=\; 23.5\ \text{min}
\quad(\text{RMS scatter } 5\%,\ \text{range } 21.5\text{--}25.5)
\label{eq:collapse}
\end{equation}
across the sweep (Fig.~\ref{fig:collapse}). Equivalently, the digitized
exponent is $t\propto V^{-1.96\pm0.06}$, statistically indistinguishable
from the $E^{-2}$ scaling that $\Lambda$ requires. Three qualifications
bound this result. Because only $V$ varies across the sweep, the collapse
constrains the field exponent of Eq.~\ref{eq:lambda} and nothing else; the
$a^{3}$ and material-property content is untested by these data. The
constant carries a systematic uncertainty of a factor $\approx$2 beyond
its 5\% scatter, since $C\propto a^{3}/d^{2}$ and depends on the choice
of size-distribution moment. And the scatter itself is consistent with
digitization noise alone (a Monte Carlo with 5\% reading error in $t$ and
2\% in $V$ reproduces it), so it should be read as an upper bound on the
law's intrinsic scatter. Three further points
follow. (i)~The lowest voltage that still cleared sits at $\Lambda\approx0.26$,
and three of the seven cleared points lie below the central
$\Lambda^{*}=0.8$: the data are consistent with the a-priori threshold
\emph{band} of Eq.~\ref{eq:lambdastar} only at its lower edge, indicating
that the central near-contact amplification is too conservative and the
true amplification is nearer the top of its 2--8 range. Below corona onset,
where the field in the oil is negligible and $\Lambda$ is effectively
zero, no clearing was observed in 110\,min, consistent with the gate. (ii)~The collapse arbitrates a genuine ambiguity in the device physics.
Writing the field as $(V-V_0)/d$ and profiling the collapse quality over
the onset offset $V_0$ gives RMS spreads of 5.2, 8.9, 15.6, 33, and 58\%
at $V_0=0$, 1, 2, 4, and 6.3\,kV: the data prefer $V_0=0$ and exclude
offsets beyond $\sim$1\,kV at the level of the observed scatter,
consistent with $V/d$ approximating the average field in the oil, as the
device's current--voltage characterization also indicates (the offset test
varies $V_0$ only; it does not by itself pin the gap).
(The corona \emph{onset} of the source device remains physically real;
what the data reject is a large standing offset in the oil-field reading
above onset.)
(iii)~The collapse does \emph{not} extend to the water-cut sweep of
Ref.~\cite{panat2025}, nor should it: $\Lambda$ answers whether a pair can
merge, not how much water must move or how much light must clear; the
water-cut axis requires a second, load-like group. Two further disclosures are essential. The collapse validates the
$E^{-2}$ scaling but \emph{not} the framework's absolute kinetics across
the sweep: at the sweep's low-field points our barrier-gated model predicts
clearing times that are orders of magnitude too long (the exponential gate
$e^{A_{nc}\Lambda-\Delta G/k_BT}$ effectively closes below
$\sim$3\,kV/cm at this droplet size), whereas the experiments cleared in
11--95\,min. The model's absolute kinetics are supported only near the
high-field reference condition; the sweep data falsify the central barrier
calibration at low fields, again pointing to stronger near-contact
amplification or to field-driven transport (electrohydrodynamic
circulation) that bypasses the gate. And the collapse is licensed only
within the sweep's range: extrapolating $t=23.5\,\text{min}/\Lambda$ to
the reference batch condition predicts minutes against the $\sim$60\,min
measured there, indicating that the sweep and the reference batch differ in
effective field or geometry in a way the published information does not
fully constrain. Resolving these tensions requires instrumented-field
source data, and we flag them as the sharpest open questions the
dimensionless analysis exposes.

On the absolute number the model \emph{overpredicts}: the central band of
90--125\,min lies a factor 1.5--2.1 above the single measured value of
$\sim$60\,min, with no stated experimental uncertainty or replicates to
compare against, so the comparison establishes order-of-magnitude
correctness with a systematic slow bias, not quantitative validation. Two
mitigating observations bound its meaning. The omitted transport channels
(electrohydrodynamic stirring, documented directly by the source
experiments) act in the observed direction. And the endpoint itself is
stiff: moving the transmittance threshold by one percentage point moves the
predicted $t_{95}$ by $-27\%$/$+46\%$ (81, 110, and 161\,min at 94, 95,
and 96\%), so part of any gap at this endpoint reflects the endpoint's own
sensitivity rather than cascade physics.

\begin{figure*}[!t]
\centering
\includegraphics[width=0.9\textwidth]{figs/fig_lambda.png}
\caption{\textbf{Validation against published data.} Left: the seven
digitized voltage-sweep points of Ref.~\cite{panat2025} against
$t=23.5\,\text{min}/\Lambda$; the flow-cell demonstration (open green) lies a factor 2.5 above the
extrapolated line, and the batch reference measurement (orange cross) and
the predicted band at the same $\Lambda$ lie an order of magnitude above
it, the extrapolation tension discussed in the text. Grey band: below
corona onset, no clearing observed; dotted line: the a-priori threshold
$\Lambda^{*}$. Right: the collapse, $t\cdot\Lambda$ with 5\% RMS
scatter.}
\label{fig:collapse}
\end{figure*}

\subsection{Why brine is the stubborn phase}

A bench observation supplied as direction only, never fitted, is that saline
droplets are markedly more stable under the field than fresh water. The
framework was given the two literature mechanisms and the slowdown emerged
in every route. First, electrolytes stiffen Span-type interfacial films:
salting-out packs the sorbitan monolayer and raises its mechanical strength
\cite{opawale1998}, the same interfacial-rigidity route by which asphaltene
films stabilize crude emulsions~\cite{kilpatrick2012}, entering here as a few additional $k_BT$ of barrier
whose effect is exponential and therefore mild at 7.5\,kV/cm but decisive at
low field. Second, brine's conductivity ($\sim$4.6\,S/m) allows charge to
transfer within the coalescence bridge, so that high-field merges of grown drops abort into daughter droplets, the
field-driven analogue of classical partial coalescence
\cite{charlesmason1960,blanchette2006,aryafar2006,ristenpart2009,bird2009,
mousavichoubeh2011}; fresh water relaxes too slowly to bounce. The net
prediction, brine slower by $\sim$20\% at 7.5\,kV/cm (110 vs 90\,min) with
the penalty concentrated in the fines tail and growing toward lower fields,
is testable: the film mechanism predicts a smooth, roughly exponential
slowdown with ionic strength, while the bounce mechanism concentrates its
damage at high field and late times, and each can be isolated by a salinity
sweep at fixed field. We note that a naive electrostatic expectation
(saltier drops polarize more strongly) predicts the opposite sign and is
contradicted by the observation: above a conductivity of order
$10\,\sigma_{\mathrm{oil}}$, the dipole is geometry-limited and salinity
adds nothing to the driving force. Classical double-layer screening
\cite{derjaguinlandau1941,verweyoverbeek1948,israelachvili2011} is likewise
secondary here, since the draining film is oil rather than electrolyte.

\section{Discussion}

The results support a simple claim with practical consequences: the clearing
time of an electrostatically demulsified nanoemulsion is predictable from
setup information alone, to engineering accuracy, and is organized by one
dimensionless group. The two-clocks decomposition explains a long-standing
frustration, namely why single-droplet intuition fails to anticipate
macroscopic clearing: the operational endpoint is governed by the fines
tail, an emergent population property. The $\Lambda$ framework converts
device design questions into geometry on the feasibility map: droplet size
sets the required field through $E_{\mathrm{crit}}\!\propto\!a^{-3/2}$, and
the measured collapse $t\!\approx\!23.5\,\text{min}/\Lambda$ prices the
speed of any operating point for this emulsion class.

Three limitations bound the claims. Validation rests on one experimental
system (seven voltage points, one absolute batch time, one flow-cell time);
the near-contact amplification $A_{nc}$ is an order-one parameter carried as
a band rather than computed, and it sets the precise location of the
feasibility cliff; and convective transport (electrohydrodynamic
circulation, shear in flow cells) is omitted, plausibly explaining both the
factor-of-two slow bias and the flow cell's additional speed. Each
limitation is also a concrete measurement: dynamic light scattering on the
residual haze after clearing would locate the fines cliff and hence
$A_{nc}$; a salinity sweep at fixed field would separate the two brine
mechanisms; and the learning agent's discovery that, under strong partial
coalescence, a field \emph{schedule} (ignite hard, then throttle to stop
bouncing) beats any constant field is directly testable with a programmable
supply.

A further caveat: the three routes share the same pair-physics kernel
layer, so they are three population-integration schemes over a common,
separately verified closure set; their mutual agreement (21\%) guards
against discretization and implementation error, which is exactly what the
development history showed it catching, and not against a shared wrong
model. Only the comparison with measurement tests the physics. More broadly, the no-calibration, triangulated protocol used here, three solvers
with disjoint failure modes, predictions frozen before data contact, and a
public verification suite, is inexpensive discipline for soft-matter
kinetics modeling generally, where post-hoc calibration is endemic and
plausible-but-wrong models are cheap to produce.

\section{Materials and Methods}

\subsection{Setup parameters}
All inputs are setup facts from Ref.~\cite{panat2025} or literature property
values: hexadecane $\mu=3.03$\,mPa\,s, $\rho_o=773$\,kg/m$^3$,
$\varepsilon_{\mathrm{oil}}=2.05$, $n_o=1.434$; 0.5\,M brine
$\rho_w=1017$\,kg/m$^3$, $n_w=1.338$, $\sigma_w=4.6$\,S/m; interfacial
tension with Span~80, $\gamma=4$\,mN/m; oil conductivity
$2\times10^{-8}$\,S/m; cell depth $H=3.37$\,cm; optical path $L=1$\,cm;
$T=298$\,K. Initial condition: lognormal, median diameter 175\,nm,
$\sigma_g=1.35$, $\phi=0.02$. The field is delivered by unipolar corona
space-charge injection across the air gap~\cite{sigmond1982,panat2025}.
Drops are treated as spherical and non-deforming, valid since the operating
field is well below the Taylor instability limit for these sizes
\cite{taylor1964,allanmason1962,berg2010}. Shear-driven collision kernels
\cite{saffmanturner1956} and hindered-settling corrections
\cite{richardsonzaki1954} are neglected, consistent with a quiescent batch
at $\phi=2\%$; both would matter in flow systems or at high water cut.

\subsection{Pair kernels}
Aligned-pair contact energy
$U_{ij}=8\pi\varepsilon_0\varepsilon_{\mathrm{oil}}K^2E^2a_i^3a_j^3/(a_i+a_j)^3$,
reducing to $\Lambda k_BT$ for equal spheres (Eq.~\ref{eq:lambda}); for
$a_j\!\gg\!a_i$, $U\!\to\!8\Lambda_i k_BT$ (the scavenging rule). Transport:
Brownian kernel $K_B=(2k_BT/3\mu)(a_i+a_j)(a_i^{-1}+a_j^{-1})$ enhanced by
the capture factor $\max(1,1.12\,\Lambda_{ij}^{1/3})$~\cite{fuchs1934}, plus
sweep coagulation $K_g=\pi R_c^2|\Delta v_{\mathrm{St}}|$ with capture
radius $R_c=\min\!\big(\Lambda_{ij}^{1/3}(a_i{+}a_j),\,(10A/\Delta
v)^{1/4}\big)$, where $A=4\varepsilon_0\varepsilon_{\mathrm{oil}}K^2E^2
a_i^3a_j^3/(\mu\,a_{\min})$ is the dipole-drift coefficient; the second
bound enforces that a partner must drift in within the transit time
\cite{pruppacher1997}. Coalescence probability per encounter:
$\min\{1,\exp(A_{nc}\Lambda_{ij}-\Delta G/k_BT)\}$ with $A_{nc}=4$ central
(2--8 band)~\cite{davis1964}; sub-unity coalescence efficiencies of this
kind are standard in film-drainage-limited collision modeling
\cite{chesters1991,ivanov1999,braziersmith1972,charlesmason1960}. Partial coalescence: probability
$g_\sigma\,\mathrm{Bo}^2/(\mathrm{Bo}^2+\mathrm{Bo}_c^2)$ with electric Bond
number $\mathrm{Bo}=\varepsilon_0\varepsilon_{\mathrm{oil}}(3E)^2
a_{\min}/\gamma$ and charge-supply gate
$g_\sigma=\tau_b/(\tau_b+\tau_e)$. This functional form, the critical
value $\mathrm{Bo}_c=0.1$, the local-field factor 3, the $+4\,k_BT$ brine
barrier increment, and the daughter recycling rule (half the smaller
drop's volume at half its radius) are our model constructions, motivated by
but not derived from Refs.~\cite{bird2009,ristenpart2009,
mousavichoubeh2011,opawale1998}; they carry order-one uncertainty and are
exercised in the sensitivity bands. Settling:
Stokes, $v_{\mathrm{St}}=2\Delta\rho g a^2/9\mu$; removal rate $v/H$.

\subsection{Route 1: analytic cascade}
Ignition frequency from the lognormal-averaged kernel; growth as an explicit
octave ladder, $t_{\mathrm{growth}}=\sum_k\tau_{\mathrm{oct}}(a_k)$ with
$\tau_{\mathrm{oct}}=2/(K_{\mathrm{eff}}n)$ and $n=3\phi/4\pi a^3$;
gravity endgame $t_{\mathrm{settle}}=(H/2)/v_{\mathrm{St}}(a_h)$ with
handoff radius $a_h$ minimizing the sum. The ladder is top-heavy in the
Brownian regime and bottom-heavy in the sweep regime, so no single-rung
approximation is valid.

\subsection{Route 2: population balance}
Sectional Smoluchowski equations~\cite{ramkrishna2000} on 123 geometric
bins (25\,nm--300\,\textmu m),
fixed-pivot deposition~\cite{kumar1996} for comparable-size merges; accretion
(mass ratio $<10^{-2}$) treated as continuous growth via upwind advection in
size space with number conservation; overflow beyond the top bin routed to
the resolved pool. Explicit conservative time stepping with
$\Delta t\le0.2/\max(\text{per-particle rate})$; the global mass audit
closes at 0.99, and refining the grid from 123 to 246 bins shifts $t_{95}$
by $+3.8\%$, so discretization error is small against the reported bands
(Fig.~\ref{fig:supp}c). Endpoint stiffness is reported in Results and
Fig.~\ref{fig:supp}b; the barrier band and the onset-offset profile are
shown in Figs.~\ref{fig:supp}a and \ref{fig:supp}d.
Endpoint: $T_{\mathrm{rel}}=\exp[-L\sum_in_iC_{\mathrm{sca}}(a_i)]$ with
Mie cross sections (relative index 0.933)~\cite{bohren1983}.

\subsection{Route 3: super-droplet Monte Carlo and learning agent}
512 computational droplets with multiplicities; pair sampling with
multiple-coalescence correction~\cite{shima2009}; independent
anomalous-diffraction optics~\cite{vandehulst1957}. The agent observes
(turbidity bin, mean-size bin, resolved-layer quartile), selects among six
field levels each 30\,s (arc risk above 8\,kV/cm), and is trained by
every-visit Monte Carlo control with constant step
size~\cite{sutton2018} (800 episodes, $\epsilon$-greedy, undiscounted);
in our experiments one-step Q-learning converged poorly here because the
near-self-transition ``field off'' action remains spuriously competitive
under bootstrapping. Predictions: the greedy-rollout distribution and
$-\max_aQ(s_0,a)$, the latter carrying a maximization bias toward
optimistic times that we note explicitly.

\subsection{Verification and data}
A twelve-check suite verifies dimensional arithmetic against raw constants,
Mie limits (Rayleigh and geometric), Stokes values, initial opacity against
the reported 0\% transmittance, zero-field stability, the collapse
statistics, the field-reading discrimination, and solver independence. The
seven voltage-sweep points were digitized from Fig.~4C of
Ref.~\cite{panat2025} (gap 5.5\,cm per its Methods; field read as $V/d$ per
its Fig.~4B). All code, parameters, and figure-generation scripts:
\url{https://github.com/sreedath/emulsion-clock}.

\appendix
\renewcommand{\thefigure}{S\arabic{figure}}
\setcounter{figure}{0}
\section*{Appendix: supplementary figures}

\begin{figure*}[!t]
\centering
\includegraphics[width=0.88\textwidth]{figs/fig_supp.png}
\caption{\textbf{Robustness of the headline numbers.} (a)~Sensitivity of
the predicted optical clearing time to the film barrier over its
20--30\,$k_BT$ band, both water phases. (b)~Stiffness of the operational
endpoint: a one-percentage-point change in the transmittance threshold
moves $t_{95}$ by $-27\%$/$+46\%$. (c)~Grid convergence of the population
balance: doubling the size resolution shifts $t_{95}$ by $+3.8\%$.
(d)~Onset-offset profile of the data collapse: the RMS spread of
$t\cdot\Lambda$ across the digitized sweep degrades monotonically with an
assumed corona offset $V_0$ in the field reading, preferring $V_0=0$.}
\label{fig:supp}
\end{figure*}

\begin{thebibliography}{99}\small
\bibitem{eow2001} J.~S. Eow, M.~Ghadiri, A.~O. Sharif, T.~J. Williams,
Electrostatic enhancement of coalescence of water droplets in oil: a review
of the current understanding. \emph{Chem. Eng. J.} \textbf{84}, 173--192
(2001).
\bibitem{tadros2009} T.~F. Tadros, \emph{Emulsion Science and Technology}
(Wiley-VCH, Weinheim, 2009).
\bibitem{atten1993} P.~Atten, Electrocoalescence of water droplets in an
insulating liquid. \emph{J. Electrostat.} \textbf{30}, 259--269 (1993).
\bibitem{panat2025} S.~Panat \emph{et al.}, Non-Laplacian air-gap
electrostatics for high-field oil--water nanoemulsion separation.
\emph{Sci. Adv.} \textbf{11}, eadz6233 (2025).
\bibitem{zhang1995} X.~Zhang, O.~A. Basaran, R.~M. Wham, Theoretical
prediction of electric field-enhanced coalescence of spherical drops.
\emph{AIChE J.} \textbf{41}, 1629--1639 (1995).
\bibitem{jones1995} T.~B. Jones, \emph{Electromechanics of Particles}
(Cambridge Univ. Press, Cambridge, 1995).
\bibitem{halsey1990} T.~C. Halsey, W.~Toor, Structure of electrorheological
fluids. \emph{Phys. Rev. Lett.} \textbf{65}, 2820--2823 (1990).
\bibitem{opawale1998} F.~O. Opawale, D.~J. Burgess, Influence of interfacial
properties of lipophilic surfactants on water-in-oil emulsion stability.
\emph{J. Colloid Interface Sci.} \textbf{197}, 142--150 (1998).
\bibitem{fuchs1934} N.~Fuchs, \"Uber die Stabilit\"at und Aufladung der
Aerosole. \emph{Z. Phys.} \textbf{89}, 736--743 (1934).
\bibitem{friedlander2000} S.~K. Friedlander, \emph{Smoke, Dust, and Haze:
Fundamentals of Aerosol Dynamics}, 2nd ed. (Oxford Univ. Press, New York,
2000).
\bibitem{davis1964} M.~H. Davis, Two charged spherical conductors in a
uniform electric field: forces and field strength. \emph{Q. J. Mech. Appl.
Math.} \textbf{17}, 499--511 (1964).
\bibitem{pruppacher1997} H.~R. Pruppacher, J.~D. Klett, \emph{Microphysics
of Clouds and Precipitation}, 2nd ed. (Kluwer, Dordrecht, 1997).
\bibitem{ristenpart2009} W.~D. Ristenpart, J.~C. Bird, A.~Belmonte,
F.~Dollar, H.~A. Stone, Non-coalescence of oppositely charged drops.
\emph{Nature} \textbf{461}, 377--380 (2009).
\bibitem{bird2009} J.~C. Bird, W.~D. Ristenpart, A.~Belmonte, H.~A. Stone,
Critical angle for electrically driven coalescence of two conical droplets.
\emph{Phys. Rev. Lett.} \textbf{103}, 164502 (2009).
\bibitem{mousavichoubeh2011} M.~Mousavichoubeh, M.~Ghadiri,
M.~Shariaty-Niassar, Electro-coalescence of an aqueous droplet at an
oil--water interface. \emph{Chem. Eng. Process.} \textbf{50}, 338--344
(2011).
\bibitem{sharma2023} E.~Sharma, R.~M. Thaokar, Population balance
modelling of an electrocoalescer. \emph{Chem. Eng. Sci.} \textbf{282},
119276 (2023).
\bibitem{anand2019} V.~Anand, R.~Patel, V.~M. Naik, V.~A. Juvekar,
R.~M. Thaokar, Modelling and particle based simulation of
electro-coalescence of a water-in-oil emulsion. \emph{Comput. Chem. Eng.}
\textbf{121}, 608--623 (2019).
\bibitem{su2023} X.~Su, X.~Yang, Y.~Yu, F.~Gao, X.~Wang, A simulation
approach to predict electrostatic coalescence performance of water-in-oil
emulsion using population balance model. \emph{Sep. Sci. Technol.}
\textbf{58}, 538--550 (2023).
\bibitem{rahmani2023} M.~Rahmani, N.~Tahouni, M.~Abbasi,
M.~H. Panjeshahi, Modeling and performance analysis of electrostatic
desalter using combined CFD-PBM method. \emph{Chem. Eng. Res. Des.}
\textbf{194}, 121--135 (2023).
\bibitem{khajeh2020} A.~Khajehesamedini \emph{et al.}, Development of
coalescence and capture kernels for the electrocoalescence process based on
batch experiments. \emph{Ind. Eng. Chem. Res.} \textbf{59}, 1277--1297
(2020).
\bibitem{kokal2005} S.~Kokal, Crude-oil emulsions: a state-of-the-art
review. \emph{SPE Prod. Facil.} \textbf{20}, 5--13 (2005).
\bibitem{zolfaghari2016} R.~Zolfaghari, A.~Fakhru'l-Razi, L.~C. Abdullah,
S.~S.~E.~H. Elnashaie, A.~Pendashteh, Demulsification techniques of
water-in-oil and oil-in-water emulsions in petroleum industry.
\emph{Sep. Purif. Technol.} \textbf{170}, 377--407 (2016).
\bibitem{mason2006} T.~G. Mason, J.~N. Wilking, K.~Meleson, C.~B. Chang,
S.~M. Graves, Nanoemulsions: formation, structure, and physical properties.
\emph{J. Phys.: Condens. Matter} \textbf{18}, R635--R666 (2006).
\bibitem{gupta2016} A.~Gupta, H.~B. Eral, T.~A. Hatton, P.~S. Doyle,
Nanoemulsions: formation, properties and applications. \emph{Soft Matter}
\textbf{12}, 2826--2841 (2016).
\bibitem{mhatre2015} S.~Mhatre \emph{et al.}, Electrostatic phase
separation: a review. \emph{Chem. Eng. Res. Des.} \textbf{96}, 177--195
(2015).
\bibitem{less2012} S.~Less, R.~Vilagines, The electrocoalescers'
technology: advances, strengths and limitations for crude oil separation.
\emph{J. Pet. Sci. Eng.} \textbf{81}, 57--63 (2012).
\bibitem{winslow1949} W.~M. Winslow, Induced fibration of suspensions.
\emph{J. Appl. Phys.} \textbf{20}, 1137--1140 (1949).
\bibitem{parthasarathy1996} M.~Parthasarathy, D.~J. Klingenberg,
Electrorheology: mechanisms and models. \emph{Mater. Sci. Eng. R}
\textbf{17}, 57--103 (1996).
\bibitem{lathamroxburgh1966} J.~Latham, I.~W. Roxburgh, Disintegration of
pairs of water drops in an electric field. \emph{Proc. R. Soc. Lond. A}
\textbf{295}, 84--97 (1966).
\bibitem{taylor1964} G.~Taylor, Disintegration of water drops in an
electric field. \emph{Proc. R. Soc. Lond. A} \textbf{280}, 383--397 (1964).
\bibitem{allanmason1962} R.~S. Allan, S.~G. Mason, Particle behaviour in
shear and electric fields I. Deformation and burst of fluid drops.
\emph{Proc. R. Soc. Lond. A} \textbf{267}, 45--61 (1962).
\bibitem{berg2010} G.~Berg, L.~E. Lundgaard, N.~Abi-Chebel, Electrically
stressed water drops in oil. \emph{Chem. Eng. Process.} \textbf{49},
1229--1240 (2010).
\bibitem{charlesmason1960} G.~E. Charles, S.~G. Mason, The coalescence of
liquid drops with flat liquid/liquid interfaces. \emph{J. Colloid Sci.}
\textbf{15}, 236--267 (1960).
\bibitem{blanchette2006} F.~Blanchette, T.~P. Bigioni, Partial coalescence
of drops at liquid interfaces. \emph{Nat. Phys.} \textbf{2}, 254--257
(2006).
\bibitem{aryafar2006} H.~Aryafar, H.~P. Kavehpour, Drop coalescence through
planar surfaces. \emph{Phys. Fluids} \textbf{18}, 072105 (2006).
\bibitem{kilpatrick2012} P.~K. Kilpatrick, Water-in-crude oil emulsion
stabilization: review and unanswered questions. \emph{Energy Fuels}
\textbf{26}, 4017--4026 (2012).
\bibitem{israelachvili2011} J.~N. Israelachvili, \emph{Intermolecular and
Surface Forces}, 3rd ed. (Academic Press, Amsterdam, 2011).
\bibitem{derjaguinlandau1941} B.~Derjaguin, L.~Landau, Theory of the
stability of strongly charged lyophobic sols and of the adhesion of
strongly charged particles in solutions of electrolytes. \emph{Acta
Physicochim. URSS} \textbf{14}, 633--662 (1941).
\bibitem{verweyoverbeek1948} E.~J.~W. Verwey, J.~Th.~G. Overbeek,
\emph{Theory of the Stability of Lyophobic Colloids} (Elsevier, Amsterdam,
1948).
\bibitem{ivanov1999} I.~B. Ivanov, K.~D. Danov, P.~A. Kralchevsky,
Flocculation and coalescence of micron-size emulsion droplets.
\emph{Colloids Surf. A} \textbf{152}, 161--182 (1999).
\bibitem{braziersmith1972} P.~R. Brazier-Smith, S.~G. Jennings, J.~Latham,
The interaction of falling water drops: coalescence. \emph{Proc. R. Soc.
Lond. A} \textbf{326}, 393--408 (1972).
\bibitem{ramkrishna2000} D.~Ramkrishna, \emph{Population Balances: Theory
and Applications to Particulate Systems in Engineering} (Academic Press,
San Diego, 2000).
\bibitem{saffmanturner1956} P.~G. Saffman, J.~S. Turner, On the collision
of drops in turbulent clouds. \emph{J. Fluid Mech.} \textbf{1}, 16--30
(1956).
\bibitem{richardsonzaki1954} J.~F. Richardson, W.~N. Zaki, Sedimentation
and fluidisation: Part I. \emph{Trans. Inst. Chem. Eng.} \textbf{32},
35--53 (1954).
\bibitem{sigmond1982} R.~S. Sigmond, Simple approximate treatment of
unipolar space-charge-dominated coronas: the Warburg law and the saturation
current. \emph{J. Appl. Phys.} \textbf{53}, 891--898 (1982).
\bibitem{mengual1999} O.~Mengual, G.~Meunier, I.~Cayre, K.~Puech,
P.~Snabre, TURBISCAN MA 2000: multiple light scattering measurement for
concentrated emulsion and suspension instability analysis. \emph{Talanta}
\textbf{50}, 445--456 (1999).
\bibitem{reddy1981} S.~R. Reddy, H.~S. Fogler, Emulsion stability:
determination from turbidity. \emph{J. Colloid Interface Sci.} \textbf{79},
101--104 (1981).
\bibitem{smoluchowski1917} M.~von Smoluchowski, Versuch einer mathematischen
Theorie der Koagulationskinetik kolloider L\"osungen. \emph{Z. Phys. Chem.}
\textbf{92}, 129--168 (1917).
\bibitem{kumar1996} S.~Kumar, D.~Ramkrishna, On the solution of population
balance equations by discretization--I. A fixed pivot technique.
\emph{Chem. Eng. Sci.} \textbf{51}, 1311--1332 (1996).
\bibitem{bohren1983} C.~F. Bohren, D.~R. Huffman, \emph{Absorption and
Scattering of Light by Small Particles} (Wiley, New York, 1983).
\bibitem{vandehulst1957} H.~C. van de Hulst, \emph{Light Scattering by Small
Particles} (Wiley, New York, 1957).
\bibitem{shima2009} S.~Shima, K.~Kusano, A.~Kawano, T.~Sugiyama,
S.~Kawahara, The super-droplet method for the numerical simulation of
clouds and precipitation: a particle-based and probabilistic microphysics
model coupled with a non-hydrostatic model. \emph{Q. J. R. Meteorol. Soc.} \textbf{135},
1307--1320 (2009).
\bibitem{sutton2018} R.~S. Sutton, A.~G. Barto, \emph{Reinforcement
Learning: An Introduction}, 2nd ed. (MIT Press, Cambridge, MA, 2018).
\bibitem{chesters1991} A.~K. Chesters, The modelling of coalescence
processes in fluid-liquid dispersions: a review of current understanding.
\emph{Chem. Eng. Res. Des.} \textbf{69}, 259--270 (1991).
\bibitem{melcher1969} J.~R. Melcher, G.~I. Taylor, Electrohydrodynamics: a
review of the role of interfacial shear stresses. \emph{Annu. Rev. Fluid
Mech.} \textbf{1}, 111--146 (1969).
\bibitem{vrij1966} A.~Vrij, Possible mechanism for the spontaneous rupture
of thin, free liquid films. \emph{Discuss. Faraday Soc.} \textbf{42},
23--33 (1966).
\end{thebibliography}

\vspace{6pt}
\noindent\textbf{Acknowledgments.} The authors thank the Vizuara AI Labs
team for infrastructure support.
\smallskip

\noindent\textbf{Author contributions.} S.P. conceived the study, performed
the source experiments, and supplied the blind-prediction constraints;
S.P., R.D., and R.D. developed the framework, performed the computations,
and wrote the manuscript.
\smallskip

\noindent\textbf{Data and code availability.} All code, inputs, digitized
data, and figure-generation scripts, under version control with the frozen
prediction state (commit \texttt{ab555b0}), are available at
\url{https://github.com/sreedath/emulsion-clock}. An interactive companion
is available at
\url{https://demulsification-prediction.vercel.app}.

\end{document}
